\documentclass{zhvt-classic}

% 使用本地源流明体（GenRyuMin2 TW）按文件路径加载，不依赖系统字体名
% 版框（内框／界栏／书口侧栏／鱼尾）改用与版心印章同一印泥棕色：
%   印章图片 yshs.png 实色 = RGB(135,68,52) = HTML 874434
\usepackage{xcolor}
\usepackage{tikz}
% sealbrown 双注册：xcolor 供 TikZ/pgf，L3 供 cls 版框
\definecolor{sealbrown}{HTML}{874434}
\ExplSyntaxOn
\color_set:nnn { sealbrown } { HTML } { 874434 }
\ExplSyntaxOff

\zhvtset{
  grid color       = sealbrown,
  main font        = GenRyuMin2TW-SB,
  main font path   = C:/Users/Administrator/texmf/fonts/opentype/public/GenRyuMin2TW/,
jiazhu font      = GenRyuMin2TW-M,
  jiazhu font path = C:/Users/Administrator/texmf/fonts/opentype/public/GenRyuMin2TW/,
  % 页码框容纳 4 字：正常页「九十七」+前置空 1 格=4 行恰好；
  % 100+ 页「一百零一」共 5 行超出，末字「一/二/三」被框裁掉（2026-10-04 实测），故放宽到 5
  shukou page number chars = 5,
  % 书口书名框高从默认 6 字减到 4 字，把省出的 2 字高拨给 chapter 框，
  % 使 6 字书名+卷名放进 9 字高的 chapter 框后不再压线（2026-10-05）
  shukou title chars = 4,
}

\title{史记}[滑稽列传]
\author{汉·司马迁}
\maker{闲叟漫记·重製}


% ===== 索隐述赞样式（2026-11-08 四次定稿 · 愚大人亲手改定参数）=====
% ⚠ 竖排前提：自定义字体家族必须带 RawFeature={vertical:+vert}；
%   simhei.ttf 加 +vhal 会整体加载失败静默回退，只写 +vert；simfang 同理。
\newCJKfontfamily{\shuzanfont}[Path=C:/Windows/Fonts/,RawFeature={vertical:+vert}]{simhei.ttf}
\newCJKfontfamily{\shuzanbody}[Path=C:/Windows/Fonts/,RawFeature={vertical:+vert}]{simfang.ttf}
% 题头：152x30pt 圆角 6pt 棕底白字；baseline (0,2) 使整框落入界行格子不压版框线；
% 白字节点 x 不补偿、y +8.2pt（像素级墨迹实测 dx/dy≈0，L/R 留白约 20pt）。
\newcommand{\shuzanhead}[1]{%
  \begin{tikzpicture}[baseline={(0,2)}]
    \fill[sealbrown, rounded corners=6pt]
      (-76pt,-15pt) rectangle (76pt,15pt);
    \node[text=white, inner sep=0pt, anchor=center] at (0pt,8.2pt)
      {\shuzanfont\fontsize{23}{23}\selectfont #1};
  \end{tikzpicture}%
}
% ★ 修正书口卷名位置：chapter 框起始位置改回 cls 默认值，避免全局影响。
%   只在每个卷首页通过 \juan 命令里的空命令微调该页位置（安全，不会跨页污染）。
\ExplSyntaxOn
\AtBeginDocument{
  \dim_set:Nn \l__zhvt_shukou_chapter_title_start_dim
    { \l__zhvt_shukou_title_height_dim + \l__zhvt_font_size_dim * 2 }
  \dim_set:Nn \l__zhvt_shukou_page_number_start_dim
    { \l__zhvt_shukou_chapter_title_start_dim + \l__zhvt_shukou_chapter_title_height_dim }
}
\ExplSyntaxOff

\begin{document}

% ★★ 版心（书口）排布覆写：把书名放到鱼尾之下（2026-10-04）
%
%   cls 默认版心自上而下：书名（title 静态框，在鱼尾之上）→ 鱼尾 → 篇名（chapter 动态框）→ 页码
%   本模板改为：          鱼尾（位置不动）→ 书名 → 篇名 → 页码分割线 → 页码
%
%   手法：把 cls 放在 title 静态框里的书名，改放进 chapter 动态框（鱼尾之下）；
%         title 静态框留空。书名取 \l__zhvt_title_str，即 \title{...} 的第一参数，
%         所以换书时只改 \title，不必在此重复写书名。
%
%   ⚠ 三条必须遵守的写法（2026-10-04 实测）：
%     ① 汉字必须用 \hbox{...} 包住。flowfram 会把框内容装进一个宽 14.4mm（= 一个 line_sep
%        = 书口栏宽）的纵向 minipage 并按该宽度折行，而整页靠 PDF /Rotate 90 显示、
%        框内文字本身不变向 —— 不包 \hbox 时「围炉夜话」4 字会被逐字折成 4 行，
%        在版面上就成了横排 4 字并横向溢出版框。
%     ② 竖排下**一个 \hbox 就是一整列**。书名与篇名要同处一列，必须放进同一个 \hbox；
%        分成两个 \hbox 时，后一个会被当成新的一列、甩到隔壁栏去
%        （实测「序」跑到书口栏外侧 14.4mm 的正文列上）。中间的 \quad 在框内是横向空白，
%        显示出来正好是书名与篇名之间的纵向间隔。
%     ③ ★ 内容盒要写成 **与框等宽的定宽盒** `\hbox to \l__zhvt_shukou_chapter_title_height_dim{…\hfil}`。
%        实测：chapter 框是把内容盒沿纵向**居中**摆放的，内容宽度一变，书名就跟着漂：
%        自然宽度的 \hbox（书名 38.4mm）落在 x=157.2mm，加篇名后（57.6mm）落在 x=152.4mm，
%        目录页（无篇名）与正文页的书名就差半格。做成等宽盒后，无论有没有篇名，
%        内容固定从框顶起排，实测两页均落在 x=142.8mm，不再漂。
\ExplSyntaxOn
% 把书口 chapter 框（书名+篇名）起始位置下移 1 字高，避免 4 字卷名顶到鱼尾
% ⚠ cls 的 \__zhvt_calculate_dimensions: 在 \AtBeginDocument 里执行，会重算该 dim；
%   必须把覆写挂到 begin document 之后，否则会被二次覆盖（2026-10-05 实测）。
% (旧方案已移到导言区，直接修改 chapter 框起始位置)
% 卷标题命令 #1=书口篇名(鱼尾) #2=正文卷名（含卷序，卷名内空格用 ~）
%
% ★★ 目录页两项修订（2026-10-05，愚大人指出）★★
%   ① 页码起点统一对齐：
%      原状 —— 目录项写成 \par\hspace{4em}卷名\hspace{4em}页码，页码起点 = 卷名实宽 + 4em。
%      卷名实宽随字数变：「卷一~梁惠王上」6 字 +~、「卷七~离娄上」5 字 +~，
%      差 1 个字位 → 卷七~卷十 4 条的页码整体比其余 10 条高 27.21pt（一个字），看着上下不齐。
%      改法 —— 卷名补全角空格（宽 = 1 字 = 27.215pt）补足到 **7 字符位**再写进目录，
%      14 条页码起点遂全部落在同一条竖线上（起点对齐；末字随页码字数自然参差）。
%      ⚠ 只对目录里的副本 \l_tmpc_tl 补，正文卷名仍用 \l_tmpa_tl 原样 —— 正文像素不受影响。
%   ② 卷名整体上移半字位：前导 \hspace 由 4em 改 3.5em。
%      实测原状目录块在版框内偏下（天头留白 139.6pt / 地脚留白 110.7pt），
%      上移 0.5em = 13.6pt 后上下基本均分，视觉居中。
\NewDocumentCommand { \juan } { mm } {
	\clearpage
	\tl_set:Nn \l__zhvt_chapter_title_tl { #1 }
	\tl_set:Nn \l_tmpa_tl { #2 }
	\tl_set_eq:NN \l_tmpc_tl \l_tmpa_tl
	\int_step_inline:nn { 7 - \str_count:n { #2 } }
	  { \tl_put_right:Nn \l_tmpc_tl { 　 } }
	\tl_set:Nx \l_tmpb_tl { \zhnum { page } }
	\tl_put_left:Nn \l_tmpb_tl { \hspace { 4em } }
	\tl_put_left:Nn \l_tmpb_tl { \l_tmpc_tl }
	\tl_put_left:Nn \l_tmpb_tl { \par \hspace { 3.5em } }
	\addtocontents { toc } { \l_tmpb_tl }
	\skip_horizontal:n { \l__zhvt_font_size_dim * 4 }
	\l_tmpa_tl
	\par
	\thispagestyle { empty }
}
\cs_set_protected:Npn \__zhvt_set_shukou_content: {
	% ★ 鱼尾上方留白（title 静态框）嵌入闲叟钤印：
	%   整页 /Rotate 90，故此处把图用 \rotatebox{90} 转正为竖章；
	%   尺寸卡在版心条幅内（书内约 7mm 宽 × 33.8mm 高，留余量不压鱼尾）。
	\setstaticcontents * { title } {
		\vbox_to_ht:nn { \l__zhvt_baseline_skip_dim } {
			\vskip 14.2pt
			\hbox_to_wd:nn { \l__zhvt_shukou_title_height_dim } {
				\hfil
				\rotatebox [ origin = c ] { 90 } { \includegraphics [ width = 7mm , height = 33.8mm ] { yshs.png } }
				\hfil
			}
		}
	}
% ★ 版心（书口）统一改法（2026-11-16 愚大人定）：18pt＋框内居中＋书口文字同版框色
%   ① 原 27.2pt 时 7 字长篇名（如「廉颇蔺相如列传」）末字溢出 chapter 框、压住页码；
%      改 \fontsize{18pt}{18pt}，前后 \hfil 令字串在框内居中（不随篇名字数漂移）。
%   ② \vbox{\vskip 8.2pt} 定量下移，抵消缩字号后栏宽方向 8.2pt 上偏（像素实测 y 中心 543.3 = 原版）。
%      ⚠ \raisebox / \vspace* 在 flowfram 框内无效；换字号须同步换补偿：16pt→10.1、19pt→7.4。
%   ③ chapter 与 page~number 均置 \color{sealbrown}（HTML 874434），与版框／界行同色；正文仍黑。
%   ⚠ 页码字位、钤印、鱼尾、界行、正文一律未动；五部合订本同此口径。
	\setdynamiccontents * { chapter }
	{ \color{sealbrown} \vbox { \vskip 8.2pt \hbox to \l__zhvt_shukou_chapter_title_height_dim
		{ \hfil \fontsize{18pt}{18pt}\selectfont \l__zhvt_title_str \quad \l__zhvt_chapter_title_tl \hfil } } }
	\setdynamiccontents * { page~number } { \color{sealbrown} \quad \zhnum { page } }
}
\__zhvt_set_shukou_content:
\ExplSyntaxOff

\maketitle

% 单篇，无多卷目录，不渲染目录页

\juan{滑稽列传}{卷一百二十六~滑稽列传第六十六}
\hspace*{2em}孔子曰。﹁六艺于治一也。礼以节人。乐以发和。书以道事。诗以达意。易以神化。春秋以义。﹂太史公曰。天道恢恢。岂不大哉。谈言微中。亦可以解纷。\\
\hspace*{2em}淳于髡者。齐之赘婿也。长不满七尺。滑稽多辩。数使诸侯。未尝屈辱。齐威王之时喜隐。好为淫乐长夜之饮。沈湎不治。委政卿大夫。百官荒乱。诸侯并侵。国且危亡。在於旦暮。左右莫敢谏。淳于髡说之以隐曰。﹁国中有大鸟。止王之庭。三年不蜚又不鸣。不知此鸟何也。﹂王曰。﹁此鸟不飞则已。一飞冲天。不鸣则已。一鸣惊人。﹂於是乃朝诸县令长七十二人。赏一人。诛一人。奋兵而出。诸侯振惊。皆还齐侵地。威行三十六年。语在田完世家中。\\
\hspace*{2em}威王八年。楚大发兵加齐。齐王使淳于髡之赵请救兵。赍金百斤。车马十驷。淳于髡仰天大笑。冠缨索绝。王曰。﹁先生少之乎。﹂髡曰。﹁何敢。﹂王曰。﹁笑岂有说乎。﹂髡曰。﹁今者臣从东方来。见道傍有禳田者。操一豚蹄。酒一盂。祝曰。﹃瓯窭满篝。汙邪满车。五穀蕃熟。穰穰满家。﹄臣见其所持者狭而所欲者奢。故笑之。﹂於是齐威王乃益赍黄金千溢。白璧十双。车马百驷。髡辞而行。至赵。赵王与之精兵十万。革车千乘。楚闻之。夜引兵而去。\\
\hspace*{2em}威王大说。置酒后宫。召髡赐之酒。问曰。﹁先生能饮几何而醉。﹂对曰。﹁臣饮一斗亦醉。一石亦醉。﹂威王曰。﹁先生饮一斗而醉。恶能饮一石哉。其说可得闻乎。﹂髡曰。﹁赐酒大王之前。执法在傍。御史在后。髡恐惧俯伏而饮。不过一斗径醉矣。若亲有严客。髡韝鞠鯱。待酒於前。时赐馀沥。奉觞上寿。数起。饮不过二斗径醉矣。若朋友交游。久不相见。卒然相睹。欢然道故。私情相语。饮可五六斗径醉矣。若乃州闾之会。男女杂坐。行酒稽留。六博投壶。相引为曹。握手无罚。目眙不禁。前有堕珥。后有遗簪。髡窃乐此。饮可八斗而醉二参。日暮酒阑。合尊促坐。男女同席。履舄交错。杯盘狼藉。堂上烛灭。主人留髡而送客。罗襦襟解。微闻芗泽。当此之时。髡心最欢。能饮一石。故曰酒极则乱。乐极则悲。万事尽然。言不可极。极之而衰。﹂以讽谏焉。齐王曰。﹁善。﹂乃罢长夜之饮。以髡为诸侯主客。宗室置酒。髡尝在侧。\\
\hspace*{2em}其后百馀年。楚有优孟。\\
\hspace*{2em}优孟。故楚之乐人也。长八尺。多辩。常以谈笑讽谏。楚庄王之时。有所爱马。衣以文绣。置之华屋之下。席以露床。啗以枣脯。马病肥死。使群臣丧之。欲以棺椁大夫礼葬之。左右争之。以为不可。王下令曰。﹁有敢以马谏者。罪至死。﹂优孟闻之。入殿门。仰天大哭。王惊而问其故。优孟曰。﹁马者王之所爱也。以楚国堂堂之大。何求不得。而以大夫礼葬之。薄。请以人君礼葬之。﹂王曰。﹁何如。﹂对曰。﹁臣请以彫玉为棺。文梓为椁。楩枫豫章为题凑。发甲卒为穿壙。老弱负土。齐赵陪位於前。韩魏翼卫其后。庙食太牢。奉以万户之邑。诸侯闻之。皆知大王贱人而贵马也。﹂王曰。﹁寡人之过一至此乎。为之柰何。﹂优孟曰。﹁请为大王六畜葬之。以垅灶为椁。铜历为棺。赍以姜枣。荐以木兰。祭以粮稻。衣以火光。葬之於人腹肠。﹂於是王乃使以马属太官。无令天下久闻也。\\
\hspace*{2em}楚相孙叔敖知其贤人也。善待之。病且死。属其子曰。﹁我死。汝必贫困。若往见优孟。言我孙叔敖之子也。﹂居数年。其子穷困负薪。逢优孟。与言曰。﹁我。孙叔敖子也。父且死时。属我贫困往见优孟。﹂优孟曰。﹁若无远有所之。﹂即为孙叔敖衣冠。抵掌谈语。岁馀。像孙叔敖。楚王及左右不能别也。庄王置酒。优孟前为寿。庄王大惊。以为孙叔敖复生也。欲以为相。优孟曰。﹁请归与妇计之。三日而为相。﹂庄王许之。三日后。优孟复来。王曰。﹁妇言谓何。﹂孟曰。﹁妇言慎无为。楚相不足为也。如孙叔敖之为楚相。尽忠为廉以治楚。楚王得以霸。今死。其子无立锥之地。贫困负薪以自饮食。必如孙叔敖。不如自杀。﹂因歌曰。﹁山居耕田苦。难以得食。起而为吏。身贪鄙者馀财。不顾耻辱。身死家室富。又恐受赇枉法。为奸触大罪。身死而家灭。贪吏安可为也。念为廉吏。奉法守职。竟死不敢为非。廉吏安可为也。楚相孙叔敖持廉至死。方今妻子穷困负薪而食。不足为也。﹂於是庄王谢优孟。乃召孙叔敖子。封之寝丘四百户。以奉其祀。后十世不绝。此知可以言时矣。\\
\hspace*{2em}其后二百馀年。秦有优旃。\\
\hspace*{2em}优旃者。秦倡侏儒也。善为笑言。然合於大道。秦始皇时。置酒而天雨。陛楯者皆沾寒。优旃见而哀之。谓之曰。﹁汝欲休乎。﹂陛楯者皆曰。﹁幸甚。﹂优旃曰。﹁我即呼汝。汝疾应曰诺。﹂居有顷。殿上上寿呼万岁。优旃临槛大呼曰。﹁陛楯郎。﹂郎曰。﹁诺。﹂优旃曰。﹁汝虽长。何益。幸雨立。我虽短也。幸休居。﹂於是始皇使陛楯者得半相代。\\
\hspace*{2em}始皇尝议欲大苑囿。东至函谷关。西至雍。陈仓。优旃曰。﹁善。多纵禽兽於其中。寇从东方来。令麋鹿触之足矣。﹂始皇以故辍止。\\
\hspace*{2em}二世立。又欲漆其城。优旃曰。﹁善。主上虽无言。臣固将请之。漆城虽於百姓愁费。然佳哉。漆城荡荡。寇来不能上。即欲就之。易为漆耳。顾难为廕室。﹂於是二世笑之。以其故止。居无何。二世杀死。优旃归汉。数年而卒。\\
\hspace*{2em}太史公曰。淳于髡仰天大笑。齐威王横行。优孟摇头而歌。负薪者以封。优旃临槛疾呼。陛楯得以半更。岂不亦伟哉。褚先生曰。臣幸得以经术为郎。而好读外家传语。窃不逊让。复作故事滑稽之语六章。编之於左。可以览观扬意。以示后世好事者读之。以游心骇耳。以附益上方太史公之三章。\\
\hspace*{2em}武帝时有所幸倡郭舍人者。发言陈辞虽不合大道。然令人主和说。武帝少时。东武侯母常养帝。帝壮时。号之曰﹁大乳母﹂。率一月再朝。朝奏入。有诏使幸臣马游卿以帛五十匹赐乳母。又奉饮飧养乳母。乳母上书曰。﹁某所有公田。原得假倩之。﹂帝曰。﹁乳母欲得之乎。﹂以赐乳母。乳母所言。未尝不听。有诏得令乳母乘车行驰道中。当此之时。公卿大臣皆敬重乳母。乳母家子孙奴从者横暴长安中。当道掣顿人车马。夺人衣服。闻於中。不忍致之法。有司请徙乳母家室。处之於边。奏可。乳母当入至前。面见辞。乳母先见郭舍人。为下泣。舍人曰。﹁即入见辞去。疾步数还顾。﹂乳母如其言。谢去。疾步数还顾。郭舍人疾言骂之曰。﹁咄。老女子。何不疾行。陛下已壮矣。宁尚须汝乳而活邪。尚何还顾。﹂於是人主怜焉悲之。乃下诏止无徙乳母。罚谪谮之者。\\
\hspace*{2em}武帝时。齐人有东方生名朔。以好古传书。爱经术。多所博观外家之语。朔初入长安。至公车上书。凡用三千奏牍。公车令两人共持举其书。仅然能胜之。人主从上方读之。止。辄乙其处。读之二月乃尽。诏拜以为郎。常在侧侍中。数召至前谈语。人主未尝不说也。时诏赐之食於前。饭已。尽怀其馀肉持去。衣尽汙。数赐缣帛。檐揭而去。徒用所赐钱帛。取少妇於长安中好女。率取妇一岁所者即弃去。更取妇。所赐钱财尽索之於女子。人主左右诸郎半呼之﹁狂人﹂。人主闻之。曰。﹁令朔在事无为是行者。若等安能及之哉。﹂朔任其子为郎。又为侍谒者。常持节出使。朔行殿中。郎谓之曰。﹁人皆以先生为狂。﹂朔曰。﹁如朔等。所谓避世於朝廷间者也。古之人。乃避世於深山中。﹂时坐席中。酒酣。据地歌曰。﹁陆沈於俗。避世金马门。宫殿中可以避世全身。何必深山之中。蒿庐之下。﹂金马门者。宦署门也。门傍有铜马。故谓之曰﹁金马门﹂。\\
\hspace*{2em}时会聚宫下博士诸先生与论议。共难之曰。﹁苏秦。张仪一当万乘之主。而都卿相之位。泽及后世。今子大夫修先王之术。慕圣人之义。讽诵诗书百家之言。不可胜数。著於竹帛。自以为海内无双。即可谓博闻辩智矣。然悉力尽忠以事圣帝。旷日持久。积数十年。官不过侍郎。位不过执戟。意者尚有遗行邪。其故何也。﹂东方生曰。﹁是固非子所能备也。彼一时也。此一时也。岂可同哉。夫张仪。苏秦之时。周室大坏。诸侯不朝。力政争权。相禽以兵。并为十二国。未有雌雄。得士者彊。失士者亡。故说听行通。身处尊位。泽及后世。子孙长荣。今非然也。圣帝在上。德流天下。诸侯宾服。威振四夷。连四海之外以为席。安於覆盂。天下平均。合为一家。动发举事。犹如运之掌中。贤与不肖。何以异哉。方今以天下之大。士民之众。竭精驰说。并进辐凑者。不可胜数。悉力慕义。困於衣食。或失门户。使张仪。苏秦与仆并生於今之世。曾不能得掌故。安敢望常侍侍郎乎。传曰。﹃天下无害菑。虽有圣人。无所施其才。上下和同。虽有贤者。无所立功。﹄故曰时异则事异。虽然。安可以不务修身乎。诗曰。﹃鼓锺于宫。声闻于外。鹤鸣九皋。声闻于天。﹄。苟能修身。何患不荣。太公躬行仁义七十二年。逢文王。得行其说。封於齐。七百岁而不绝。此士之所以日夜孜孜。修学行道。不敢止也。今世之处士。时虽不用。崛然独立。塊然独处。上观许由。下察接舆。策同范蠡。忠合子胥。天下和平。与义相扶。寡偶少徒。固其常也。子何疑於余哉。﹂於是诸先生默然无以应也。\\
\hspace*{2em}建章宫后閤重栎中有物出焉。其状似麋。以闻。武帝往临视之。问左右群臣习事通经术者。莫能知。诏东方朔视之。朔曰。﹁臣知之。原赐美酒粱饭大飧臣。臣乃言。﹂诏曰。﹁可。﹂已又曰。﹁某所有公田鱼池蒲苇数顷。陛下以赐臣。臣朔乃言。﹂诏曰。﹁可。﹂於是朔乃肯言。曰。﹁所谓驺牙者也。远方当来归义。而驺牙先见。其齿前后若一。齐等无牙。故谓之驺牙。﹂其后一岁所。匈奴混邪王果将十万众来降汉。乃复赐东方生钱财甚多。\\
\hspace*{2em}至老。朔且死时。谏曰。﹁诗云﹃营营青蝇。止于蕃。恺悌君子。无信谗言。谗言罔极。交乱四国﹄。原陛下远巧佞。退谗言。﹂帝曰。﹁今顾东方朔多善言。﹂怪之。居无几何。朔果病死。传曰。﹁鸟之将死。其鸣也哀。人之将死。其言也善。﹂此之谓也。\\
\hspace*{2em}武帝时。大将军卫青者。卫后兄也。封为长平侯。从军击匈奴。至余吾水上而还。斩首捕虏。有功来归。诏赐金千斤。将军出宫门。齐人东郭先生以方士待诏公车。当道遮卫将军车。拜谒曰。﹁原白事。﹂将军止车前。东郭先生旁车言曰。﹁王夫人新得幸於上。家贫。今将军得金千斤。诚以其半赐王夫人之亲。人主闻之必喜。此所谓奇策便计也。﹂卫将军谢之曰。﹁先生幸告之以便计。请奉教。﹂於是卫将军乃以五百金为王夫人之亲寿。王夫人以闻武帝。帝曰。﹁大将军不知为此。﹂问之安所受计策。对曰。﹁受之待诏者东郭先生。﹂诏召东郭先生。拜以为郡都尉。东郭先生久待诏公车。贫困饥寒。衣敝。履不完。行雪中。履有上无下。足尽践地。道中人笑之。东郭先生应之曰。﹁谁能履行雪中。令人视之。其上履也。其履下处乃似人足者乎。﹂及其拜为二千石。佩青緺出宫门。行谢主人。故所以同官待诏者。等比祖道於都门外。荣华道路。立名当世。此所谓衣褐怀宝者也。当其贫困时。人莫省视。至其贵也。乃争附之。谚曰。﹁相马失之瘦。相士失之贫。﹂其此之谓邪。\\
\hspace*{2em}王夫人病甚。人主至自往问之曰。﹁子当为王。欲安所置之。﹂对曰。﹁原居洛阳。﹂人主曰。﹁不可。洛阳有武库。敖仓。当关口。天下咽喉。自先帝以来。传不为置王。然关东国莫大於齐。可以为齐王。﹂王夫人以手击头。呼﹁幸甚﹂。王夫人死。号曰﹁齐王太后薨﹂。\\
\hspace*{2em}昔者。齐王使淳于髡献鹄於楚。出邑门。道飞其鹄。徒揭空笼。造诈成辞。往见楚王曰。﹁齐王使臣来献鹄。过於水上。不忍鹄之渴。出而饮之。去我飞亡。吾欲刺腹绞颈而死。恐人之议吾王以鸟兽之故令士自伤杀也。鹄。毛物。多相类者。吾欲买而代之。是不信而欺吾王也。欲赴佗国奔亡。痛吾两主使不通。故来服过。叩头受罪大王。﹂楚王曰。﹁善。齐王有信士若此哉。﹂厚赐之。财倍鹄在也。\\
\hspace*{2em}武帝时。徵北海太守诣行在所。有文学卒史王先生者。自请与太守俱。﹁吾有益於君﹂。君许之。诸府掾功曹白云。﹁王先生嗜酒。多言少实。恐不可与俱。﹂太守曰。﹁先生意欲行。不可逆。﹂遂与俱。行至宫下。待诏宫府门。王先生徒怀钱沽酒。与卫卒仆射饮。日醉。不视其太守。太守入跪拜。王先生谓户郎曰。﹁幸为我呼吾君至门内遥语。﹂户郎为呼太守。太守来。望见王先生。王先生曰。﹁天子即问君何以治北海令无盗贼。君对曰何哉。﹂对曰。﹁选择贤材。各任之以其能。赏异等。罚不肖。﹂王先生曰。﹁对如是。是自誉自伐功。不可也。原君对言。非臣之力。尽陛下神灵威武所变化也。﹂太守曰。﹁诺。﹂召入。至于殿下。有诏问之曰。﹁何於治北海。令盗贼不起。﹂叩头对言。﹁非臣之力。尽陛下神灵威武之所变化也。﹂武帝大笑。曰。﹁於呼。安得长者之语而称之。安所受之。﹂对曰。﹁受之文学卒史。﹂帝曰。﹁今安在。﹂对曰。﹁在宫府门外。﹂有诏召拜王先生为水衡丞。以北海太守为水衡都尉。传曰。﹁美言可以市。尊行可以加人。君子相送以言。小人相送以财。﹂\\
\hspace*{2em}魏文侯时。西门豹为鄴令。豹往到鄴。会长老。问之民所疾苦。长老曰。﹁苦为河伯娶妇。以故贫。﹂豹问其故。对曰。﹁鄴三老。廷掾常岁赋敛百姓。收取其钱得数百万。用其二三十万为河伯娶妇。与祝巫共分其馀钱持归。当其时。巫行视小家女好者。云是当为河伯妇。即娉取。洗沐之。为治新缯绮縠衣。间居斋戒。为治斋宫河上。张缇绛帷。女居其中。为具牛酒饭食。十馀日。共粉饰之。如嫁女床席。令女居其上。浮之河中。始浮。行数十里乃没。其人家有好女者。恐大巫祝为河伯取之。以故多持女远逃亡。以故城中益空无人。又困贫。所从来久远矣。民人俗语曰﹃即不为河伯娶妇。水来漂没。溺其人民﹄云。﹂西门豹曰。﹁至为河伯娶妇时。原三老。巫祝。父老送女河上。幸来告语之。吾亦往送女。﹂皆曰。﹁诺。﹂\\
\hspace*{2em}至其时。西门豹往会之河上。三老。官属。豪长者。里父老皆会。以人民往观之者三二千人。其巫。老女子也。已年七十。从弟子女十人所。皆衣缯单衣。立大巫后。西门豹曰。﹁呼河伯妇来。视其好丑。﹂即将女出帷中。来至前。豹视之。顾谓三老。巫祝。父老曰。﹁是女子不好。烦大巫妪为入报河伯。得更求好女。后日送之。﹂即使吏卒共抱大巫妪投之河中。有顷。曰。﹁巫妪何久也。弟子趣之。﹂复以弟子一人投河中。有顷。曰。﹁弟子何久也。复使一人趣之。﹂复投一弟子河中。凡投三弟子。西门豹曰。﹁巫妪弟子是女子也。不能白事。烦三老为入白之。﹂复投三老河中。西门豹簪笔磬折。乡河立待良久。长老。吏傍观者皆惊恐。西门豹顾曰。﹁巫妪。三老不来还。柰之何。﹂欲复使廷掾与豪长者一人入趣之。皆叩头。叩头且破。额血流地。色如死灰。西门豹曰。﹁诺。且留待之须臾。﹂须臾。豹曰。﹁廷掾起矣。状河伯留客之久。若皆罢去归矣。﹂鄴吏民大惊恐。从是以后。不敢复言为河伯娶妇。\\
\hspace*{2em}西门豹即发民凿十二渠。引河水灌民田。田皆溉。当其时。民治渠少烦苦。不欲也。豹曰。﹁民可以乐成。不可与虑始。今父老子弟虽患苦我。然百岁后期令父老子孙思我言。﹂至今皆得水利。民人以给足富。十二渠经绝驰道。到汉之立。而长吏以为十二渠桥绝驰道。相比近。不可。欲合渠水。且至驰道合三渠为一桥。鄴民人父老不肯听长吏。以为西门君所为也。贤君之法式不可更也。长吏终听置之。故西门豹为鄴令。名闻天下。泽流后世。无绝已时。几可谓非贤大夫哉。\\
\hspace*{2em}传曰。﹁子产治郑。民不能欺。子贱治单父。民不忍欺。西门豹治鄴。民不敢欺。﹂三子之才能谁最贤哉。辨治者当能别之。\\

% ===== 索隐述赞（司马贞《史记索隐》）· 统一样式 =====
% 题头「索隐述赞曰」五字并入反白框；正文每四句一列，定宽盒 \hfil 居列中；
% \fontsize{24}{41} 列距对齐全书界行（121+41k）；\clearpage 使述赞独占一页。
\clearpage
\shuzanhead{索隐述赞曰}\linebreak
{\shuzanbody\fontsize{24}{41}\selectfont
\hbox to \hsize{\hfil 滑稽鸱夷。如脂如韦。敏捷之变。学不失词。\hfil}\linebreak
\hbox to \hsize{\hfil 淳于索绝。赵国兴师。楚优拒相。寝丘获祠。\hfil}\linebreak
\hbox to \hsize{\hfil 伟哉方朔。三章纪之。\hfil}%
\par}
\end{document}
